\section{Introduction}
\label{sec:tpc39-introduction}

The orbit--sliced program in TPC--34--TPC--38 reduces a prime--pair
remainder to a coefficient--specific Hilbert energy for reflected
M\"obius divisor sums
\citep{WangTPC34,WangTPC35,WangTPC36,WangTPC37,WangTPC38}.
After exact full/proper recombination, three orbit--scale CRT coordinates
cancel all nonjoint determinant ghosts.  The remaining
congruence completion is
\begin{equation}
 \mathscr Y_{\rm alias}
 =\sum_{k=-L}^{L}\mathscr Z_k,
 \qquad
 L=\left\lfloor\frac BM\right\rfloor,
 \label{eq:tpc39-intro-alias-sum}
\end{equation}
where \(M\) is the three--modulus product, \(|F|\le B\) is the
determinant window, and \(\mathscr Z_k\) is the raw column on
\(F=kM\).  The physical equality output is only
\begin{equation}
 \mathscr Y_{\rm phys}=\mathscr Z_0.
 \label{eq:tpc39-intro-physical-zero}
\end{equation}

TPC--38 proved that \eqref{eq:tpc39-intro-alias-sum} does not control
\eqref{eq:tpc39-intro-physical-zero} in the unrestricted Hilbert column
class.  The obstruction is elementary and sharp: the columns
\(\mathscr Z_0=v\) and \(\mathscr Z_1=-v\) make the completed sum vanish
while leaving the equality column arbitrary.  Multiplicative character
refinement cannot separate them, since every joint alias satisfies
\begin{equation}
 (N+kM)N^{-1}\equiv1\pmod M.
 \label{eq:tpc39-intro-ratio-one}
\end{equation}

The present paper uses the additive quotient that remains after
\eqref{eq:tpc39-intro-ratio-one}.  Put
\begin{equation}
 q=L+1,
 \qquad \omega=\ee{1/q}.
 \label{eq:tpc39-intro-minimal-q}
\end{equation}
Then the vectors
\begin{equation}
 \mathscr Y_\nu
 =\sum_{k=-L}^{L}\omega^{\nu k}\mathscr Z_k,
 \qquad 0\le\nu<q,
 \label{eq:tpc39-intro-twisted-bank}
\end{equation}
obey
\begin{equation}
 \boxed{
 \mathscr Z_0=\frac1q\sum_{\nu=0}^{q-1}\mathscr Y_\nu.}
 \label{eq:tpc39-intro-exact-recovery}
\end{equation}
This is not a full discrete Fourier inversion: there are only \(q=L+1\)
measurements for \(2L+1\) alias columns.  Positive and negative aliases
are folded together modulo \(q\), but the zero residue class contains
only \(k=0\).  This is exactly the amount of information needed for
equality recovery.

\subsection{The finite quotient projector}

The physical form of \eqref{eq:tpc39-intro-exact-recovery} is the identity
\begin{equation}
 \boxed{
 \one_{M\mid F}\frac1q\sum_{\nu=0}^{q-1}
 \ee{\nu F/(qM)}
 =\one_{F=0}
 \qquad(|F|\le B).}
 \label{eq:tpc39-intro-quotient-projector}
\end{equation}
Indeed, the left side is \(\one_{qM\mid F}\), and
\(qM>B\).  The factor \(\one_{M\mid F}\) is supplied by the exact
full/proper recombination
\begin{equation}
 C_N(F)+P_N(F)=\one_{M\mid F}
 \label{eq:tpc39-intro-face-recombination}
\end{equation}
proved in TPC--38.  It cannot be omitted: the additive average alone
is generally nonzero when \(M\nmid F\).  Only after full/proper
recombination restricts the determinant to \(F=kM\) does the quotient
average select the multiples of \(qM\).

Thus \eqref{eq:tpc39-intro-quotient-projector} is an additive projector on
the quotient
\begin{equation}
 M\Z/qM\Z\cong\Z/q\Z.
 \label{eq:tpc39-intro-quotient-group}
\end{equation}
It converts the congruence completion into exact equality without adding
a fourth multiplicative character coordinate.

\subsection{Folded Parseval and the exact spectrum}

For Hilbert--valued columns, character orthogonality on
\eqref{eq:tpc39-intro-quotient-group} gives
\begin{equation}
 \boxed{
 \frac1q\sum_{\nu=0}^{q-1}\|\mathscr Y_\nu\|^2
 =\|\mathscr Z_0\|^2+
 \sum_{r=1}^{L}\|\mathscr Z_r+\mathscr Z_{r-q}\|^2.}
 \label{eq:tpc39-intro-folded-Parseval}
\end{equation}
The associated normalized tomography Gram has spectrum
\begin{equation}
 2\quad(L\text{ times}),
 \qquad
 1\quad(1\text{ time}),
 \qquad
 0\quad(L\text{ times}).
 \label{eq:tpc39-intro-spectrum}
\end{equation}
The eigenvalue--one direction is the zero alias.  The kernel consists of
antisymmetric pairs
\begin{equation}
 \mathscr Z_r=-\mathscr Z_{r-q},
 \qquad 1\le r\le L,
 \label{eq:tpc39-intro-kernel}
\end{equation}
which do not affect \(\mathscr Z_0\).  Hence this bank is strictly weaker
than a full alias trace and much weaker than an alias--Gram operator bound,
but it is sufficient for the required column.

\subsection{Minimal sampling and optimal recovery amplification}

The sample count in \eqref{eq:tpc39-intro-minimal-q} is best possible.
Suppose phase nodes \(z_1,\ldots,z_m\in\C^\times\) and complex recovery
weights are required to recover the central coefficient of every Laurent
polynomial supported on \([-L,L]\).  If \(m\le L\), the polynomial
\begin{equation}
 Q(z)=\prod_{j=1}^{m}(z-z_j)
 \label{eq:tpc39-intro-annihilator}
\end{equation}
has degree at most \(L\), vanishes at every measurement node, and has
nonzero constant coefficient.  Exact recovery is therefore impossible.
The regular \((L+1)\)-gon attains the lower bound.

If \(a\) is the recovery vector for \(m\) measurements, its normalized
noise amplification is
\begin{equation}
 \cA=m\|a\|_2^2.
 \label{eq:tpc39-intro-amplification}
\end{equation}
The constant moment forces \(\cA\ge1\).  Equal weights on the regular
polygon give \(\cA=1\).  Under positive weights, the minimally sampled
rule is rigid: the weights are equal and the nodes are a rotation of the
regular polygon.

\subsection{The new physical gate}

For the physical determinant \(F=su-N\), define the twisted face fields
\begin{equation}
 C_{\nu,N}(F)=\ee{\nu F/(qM)}C_N(F),
 \qquad
 P_{\nu,N}(F)=\ee{\nu F/(qM)}P_N(F).
 \label{eq:tpc39-intro-twisted-fields}
\end{equation}
Their coefficient masses are exactly those of the untwisted fields.
After recombination, their outputs are the vectors in
\eqref{eq:tpc39-intro-twisted-bank}.  Therefore the explicit conditional
input
\begin{equation}
 \frac1q\sum_{\nu=0}^{q-1}\|\mathscr Y_\nu\|^2
 \ll X^{o(1)}Q^2J(2L+1)
 \label{eq:tpc39-intro-twisted-gate}
\end{equation}
implies the physical equality gate.  At the endpoint,
\begin{equation}
 2L+1\asymp q\asymp X^{1/400+o(1)},
 \qquad
 Q^2J(2L+1)=X^{o(1)}Q^3/J.
 \label{eq:tpc39-intro-endpoint}
\end{equation}

Equation \eqref{eq:tpc39-intro-twisted-gate} is not proved here.  Folded
Parseval shows that it retains \(\|\mathscr Z_0\|^2\) as a positive
summand.  In the terminal factorization \(s=1\), the twists create a
short Fourier transform along the progression \(N+kM\), but the square
still contains the original affine four--M\"obius term and long--gap
correlations.  Additive tomography repairs the logical recovery gap; it
does not manufacture arithmetic cancellation.

\subsection{Organization and claim boundary}

\Cref{sec:tpc39-interface} states the inherited physical interface.
\Cref{sec:tpc39-projector} proves the quotient projector and exact Hilbert
recovery.  \Cref{sec:tpc39-folded} proves folded Parseval and the complete
spectrum.  \Cref{sec:tpc39-minimality} proves sample minimality,
recovery--amplification optimality, and positive--rule rigidity.
\Cref{sec:tpc39-completion-comparison} compares phase recovery with pure
divisibility completions.  \Cref{sec:tpc39-physical-gate} formulates the
twist--stable physical gate.  \Cref{sec:tpc39-terminal} records its
terminal M\"obius content.  \Cref{sec:tpc39-partial-random} treats partial
and random banks, and \Cref{sec:tpc39-next-gate} gives the corrected route
tree.

All quotient, Fourier, spectrum, minimality, and finite divisibility
results are unconditional.  The paper proves no twist--stable physical
energy estimate, no new M\"obius autocorrelation estimate, no affine
Chowla theorem, no failure of sieve parity, no Hardy--Littlewood
prime--pair asymptotic, and no twin--prime theorem.
